\section*{Kinematyka}

    \subsection*{Opis ruchu we współrzędnych kartezjańskich}
    
        Ogólne równanie ruchu punktu materialnego po elipsie dane jest wzorami
    
        \begin{align*}
            x(t) &= A \cos \phi(t) \\
            y(t) &= B \sin \phi(t)
        \end{align*}
        
        Wiedząc, że $\phi(t) = -\frac{\pi}{4}t^2 + \frac{\pi}{2}$, A = 2, B = 1, oblicz i narysuj wektory wodzące, wektory prędkości oraz wektory przyspieszenia w chwilach t = {0; 1; 2}.
    
        \begin{table}[H]\small
        \begin{tabular}{llllllllllll} \hline
        $t$  & $\phi$ & $\sin\phi$ & $\cos\phi$ & $x$  & $y$   & $v_x$ & $v_y$ & $|v|$ & $a_x$ & $a_y$  & $|a|$ \\ \hline
        0.00 & 1.57   & 1.00  & 0.00 & 0.00 & 1.00  & 0.00  & 0.00  & 0.00  & 3.14   & 0.00   & 3.14  \\ \hline
        1.00 & 0.79   & 0.71  & 0.71 & 1.41 & 0.71  & 2.22  & -1.11 & 2.48  & -1.27  & -2.86  & 3.12  \\
        2.00 & -1.57  & -1.00 & 0.00 & 0.00 & -1.00 & -6.28 & 0.00  & 6.28  & -3.14  & 9.87   & 10.36 \\
        3.00 & -5.50  & 0.71  & 0.71 & 1.41 & 0.71  & 6.66  & -3.33 & 7.45  & -29.18 & -16.81 & 33.68 \\
        4.00 & -11.00 & 1.00  & 0.00 & 0.00 & 1.00  & 12.57 & 0.00  & 12.57 & 3.14   & -39.48 & 39.60 \\ \hline
        \end{tabular}
        \end{table}
    
    \subsection*{Opis ruchu we współrzędnych naturalnych}
    
        Ruch punktu dany jest równaniem parametrycznym
        
        \begin{align*}
            x &= 2 \sqrt{t} \\
            y &= x^2
        \end{align*}
        
        Narysuj tor punktu i oblicz całkowitą prędkość (szybkość) oraz przyspieszenie styczne i normalne w chwilach t = { 0.1; 0.2; 0.3; 0.4}. Narysuj wektory prędkości i przyspieszenia oraz ich składowe we współrzędnych kartezjańskich oraz naturalnych.
    
        \begin{table}[H] \small
        \begin{tabular}{lllllll} \hline
        czas  & x     & y     & vx    & vy    & ax      & ay    \\ \hline
        0.100 & 0.632 & 0.400 & 3.162 & 4.000 & -15.811 & 0.000 \\
        0.200 & 0.894 & 0.800 & 2.236 & 4.000 & -5.590  & 0.000 \\
        0.300 & 1.095 & 1.200 & 1.826 & 4.000 & -3.043  & 0.000 \\
        0.400 & 1.265 & 1.600 & 1.581 & 4.000 & -1.976  & 0.000 \\ \hline
        \end{tabular}
        \end{table}
        
        \begin{table}[H] \small
        \begin{tabular}{lllllll} \hline
        czas  & lenv  & evx   & evy   & lena   & at     & an     \\ \hline
        0.100 & 5.099 & 0.620 & 0.784 & 15.811 & -9.806 & 12.403 \\
        0.200 & 4.583 & 0.488 & 0.873 & 5.590  & -2.728 & 4.880  \\
        0.300 & 4.397 & 0.415 & 0.910 & 3.043  & -1.263 & 2.768  \\
        0.400 & 4.301 & 0.368 & 0.930 & 1.976  & -0.727 & 1.838  \\ \hline
        \end{tabular}
        \end{table}